\section{Agentes de aplicación frente a capacidad del modelo}

La sección anterior trató el PMF; esta sección vuelve a la cuestión competitiva central del emprendimiento en aplicaciones de IA: si los modelos son cada vez más potentes, ¿qué barreras le quedan a una empresa de aplicaciones? La respuesta de Lovart es el flujo de trabajo vertical. La ventaja coyuntural de los modelos se homogeneiza con rapidez, pero en el escenario del diseñador el contexto, el criterio estético, la gestión de activos, la iteración de versiones y el proceso de entrega siguen necesitando un producto que los sostenga.

\begin{figure}[H]
\centering
\includegraphics[width=0.94\textwidth]{agent-app-vs-model.png}
\caption{Agentes de aplicación frente a capacidad del modelo: el emprendimiento en aplicaciones debe consolidar flujo de trabajo y posicionamiento en la mente del usuario dentro de la ventana de ventaja que abren los modelos. Diagrama conceptual propio, elaborado a partir de la conversación de 00:55:00--01:10:30.}
\end{figure}

\begin{knowledgebox}{Lectura de la figura: la ventaja de los modelos no es una barrera para la aplicación}
Que el modelo base se vuelva más potente abre una ventana, pero también hace que las funciones similares se generalicen con rapidez. La empresa de aplicaciones debe convertir ese periodo de oportunidad en flujo de trabajo, posicionamiento en la mente del usuario, retorno de datos y experiencia vertical.
\end{knowledgebox}

\subsection{Los Agent generales y los Agent verticales coexistirán}

La sección anterior explicó que la barrera de una aplicación proviene del flujo de trabajo; esta sección analiza su relación con los puntos de entrada generales. En la entrevista se menciona que los Agent generales y los Agent verticales pueden coexistir. Los Agent generales se convertirán en punto de entrada, pero los Agent verticales ofrecen, en escenarios profesionales, herramientas más profundas, un contexto más sólido y una mejor experiencia. Lovart podría ser invocado por un Agent general y también convertirse en el punto de entrada especializado para el escenario de diseño.

\begin{importantbox}{Criterio para un Agent vertical}
Si un escenario tiene contexto profesional, herramientas especializadas, iteración repetida, estándares de entrega y usuarios dispuestos a pagar, vale la pena construir un Agent vertical. El diseño cumple exactamente estas condiciones.
\end{importantbox}

\begin{knowledgebox}{Modelo de coexistencia general/vertical}
\begin{tabularx}{\textwidth}{p{0.22\textwidth}p{0.34\textwidth}X}
\toprule
Forma & Función & Posición de Lovart \\
\midrule
Punto de entrada de Agent general & Recibe la intención del usuario y distribuye tareas & Podría invocar a Lovart para completar tareas de diseño. \\
Espacio de trabajo de Agent vertical & Profundiza en escenarios profesionales y cadenas de herramientas & Atiende directamente a diseñadores y equipos creativos. \\
Modelo base & Aporta capacidades de generación, comprensión y razonamiento & Lovart debe aprovecharlo, pero no puede depender solo del modelo base. \\
Ecosistema de plugins/herramientas & Conecta distintas capacidades profesionales & Lovart puede convertirse en un nodo de capacidad de diseño. \\
\bottomrule
\end{tabularx}
\end{knowledgebox}

\subsection{Modelo de negocio y riesgos de derechos de autor y de marca}

Aclarada la relación entre lo general y lo vertical, esta sección pasa a la comercialización y los riesgos. El modelo de negocio de un Agent de diseño no es sencillo. Puede cobrar mediante suscripción individual, ofrecer un espacio de trabajo para equipos o cobrar por proyecto, por volumen de material o por licencia comercial. Pero cuanto más se adentra en el trabajo de diseño real, más importan los derechos de autor, el origen del material, la coherencia de marca y la seguridad de los activos del cliente. Los diseñadores y los clientes empresariales no solo preguntan «¿puede generarlo?», sino también «¿el resultado puede usarse comercialmente?», «¿cumple con las normas de la marca?» y «¿el equipo puede seguir editándolo?».

\begin{knowledgebox}{Tabla de modelos de negocio de un Agent de diseño}
\begin{tabularx}{\textwidth}{p{0.20\textwidth}p{0.34\textwidth}X}
\toprule
Modelo & Ventaja & Riesgo \\
\midrule
Suscripción individual & Adopción rápida, crecimiento sencillo & Ticket promedio bajo; fácil de sustituir por herramientas generales. \\
Espacio de trabajo para equipos & Entra en el proceso de colaboración; mayor retención & Requiere gestión de permisos, versiones y activos. \\
Cobro por proyecto & Más ligado al valor para el cliente & Fijación de precios compleja; mayor responsabilidad de entrega. \\
Licencia comercial/material & Permite vincularse al valor de marca y de derechos de autor & Límites legales, de derechos de autor y de responsabilidad más complejos. \\
\bottomrule
\end{tabularx}
\end{knowledgebox}

\begin{importantbox}{Umbral empresarial de un Agent de diseño}
Al entrar de verdad en el proceso de diseño de una empresa, Lovart no solo debe resolver la calidad de generación, sino también la seguridad de los activos, la explicabilidad de los derechos de autor, la coherencia de marca, la colaboración en equipo y la entrega editable. Esta es la zona más exigente para un Agent vertical.
\end{importantbox}

\subsection{Estructura de tres niveles de las barreras de una aplicación}

La sección anterior explicó que la comercialización empujará a Lovart hacia flujos de trabajo más profundos; esta sección desglosa las barreras que podría formar. Si Lovart quiere resistir a largo plazo el avance de las empresas de modelos hacia la capa de aplicaciones, necesita tres niveles de barreras. El primero es la barrera de experiencia: que el diseñador lo encuentre más cómodo que una ventana de chat general. El segundo es la barrera de flujo de trabajo: los proyectos del equipo, los activos, las versiones y las entregas están todos dentro. El tercero es la barrera de ecosistema: la conexión con herramientas de diseño, bancos de material, sistemas de marca y procesos de colaboración con clientes.

\begin{knowledgebox}{Estructura de tres niveles de las barreras de una aplicación}
\begin{tabularx}{\textwidth}{p{0.20\textwidth}p{0.34\textwidth}X}
\toprule
Nivel & Cómo se forma & Riesgo de ser copiada \\
\midrule
Barrera de experiencia & Mejor interacción para el diseñador y mejores procesos predeterminados & Los productos de modelos grandes pueden alcanzarla con facilidad. \\
Barrera de flujo de trabajo & Acumulación de proyectos, equipos, versiones y activos & El costo de migración es relativamente alto. \\
Barrera de ecosistema & Conexión con herramientas, material, marca y sistemas de entrega & La más difícil de copiar, pero la más lenta de construir. \\
\bottomrule
\end{tabularx}
\end{knowledgebox}

\begin{warningbox}{El mayor riesgo de una aplicación vertical es un «flujo de trabajo delgado»}
Si la aplicación solo le pone una interfaz al modelo, los Agent generales la absorberán muy pronto. Solo cuando el flujo de trabajo, los datos y las relaciones de colaboración se consolidan de verdad, la aplicación vertical tiene una posición duradera.
\end{warningbox}

\subsection{Resumen de la sección}

El reto a largo plazo de Lovart es convertir la ventana de capacidad de los modelos en una barrera de flujo de trabajo para diseñadores. Los Agent generales concentrarán los puntos de entrada, pero los Agent verticales aún pueden tener una zona más exigente propia en los escenarios profesionales.
